\section{Models and training}
\label{si:training}

\subsection{Readout of the structure-only network}

The structure-only network S maps the pristine supercell, with the oxygen atom to be removed marked, to the
formation energy of the neutral vacancy (Figure~\ref{fig:readout}). Successive MEGNet blocks
\citep{Chen2019graph} aggregate information over growing coordination shells around every atom, starting from the
learned element embeddings, the vacancy flag, a Gaussian expansion of each interatomic distance and a global state
of the cell. The final step joins three learned vectors: the description of the vacancy site, a pooled description
of all atoms of the host (Set2Set or the mean, a tuned choice) and the final global state; a small multilayer
network maps them to one energy. The split follows the variance of the target. The pooled host vector and the
global state carry the host-level thermodynamic term, the cost of releasing an oxygen atom from that oxide, and
the vacancy-site vector carries the local environment that separates the sites of one host. In D (the D-state
injection), the host-electronic descriptors enter the initial global state and the site-electronic descriptors the
initial description of the vacancy site, each standardised with training-set statistics, so that each class enters
where its physics acts; the readout itself is unchanged. The readout is the same for both networks, the
originally registered network (v1) and the primary network (v2).

\begin{figure}[htbp]
\centering
\sifig{fig_readout}
\caption{Readout of the structure-only network S, from the initial features of the vacancy site \xvzero{}, the other atoms
\xizero{}, the bonds \eijzero{} and the cell \uzero{}, through $n$ MEGNet blocks, to the vacancy-site vector
$\mathbf{h}_\mathrm{v}$, the pooled vector $\bar{\mathbf{h}}$ and the global state $u$ that a multilayer perceptron
maps to $E_\mathrm{f}$. Green arrows mark where D (D-state) receives the host-electronic descriptors, at \uzero{},
and the site-electronic descriptors, at \xvzero{}.}
\label{fig:readout}
\end{figure}

\subsection{Training protocol, search space and tuned values}

Every network was trained on the standardised formation energy with an absolute-error loss and the AdamW
optimiser; the learning rate was halved (factor \numCfgTrainingSchedulerFactor{}) when the validation error had
not improved for \numCfgTrainingSchedulerPatience{} passes over the training hosts. The validation hosts that
decide early stopping are carved out of the training hosts of each budget (a fraction of
\numCfgTrainingValidationFrac{}, at least \numCfgTrainingValidationMinHosts{} hosts), so a budget counts every
host whose DFT label a model consumes. Training stops when the validation error has not improved for
\numCfgTrainingEarlyStoppingPatience{} passes, or at the cap of \numCfgTrainingMaxEpochs{} passes, and the weights
of the best pass are restored.

The settings of each network were tuned once, before any test host was evaluated, on the nested training sets of
resample \numCfgTuningAnchorResample{} at three anchor budgets, \numCfgTuningAnchorsItemOne{},
\numCfgTuningAnchorsItemTwo{} and \numCfgTuningAnchorsItemThree{} hosts (\texttt{docs/tuning.md}). Each model and
anchor had one study of \numCfgTuningOptunaTrialsPerModelPerAnchor{} complete trials with a tree-structured
Parzen sampler, scored by the best validation error; at the largest anchor, \numTuningSummaryStudiesSASixHundredFiftyFourValidationNTrainHosts{}
hosts were fitted and \numTuningSummaryStudiesSASixHundredFiftyFourValidationNValHosts{} validated. The search
space was the same for every model (Table~\ref{tab:space}). A training budget uses the values of its anchor: the
budgets \numCfgBudgetsHostsItemOne{} and \numCfgBudgetsHostsItemTwo{} those of anchor
\numCfgTuningBudgetToAnchorTwentyFive{}, the budgets \numCfgBudgetsHostsItemThree{} and
\numCfgBudgetsHostsItemFour{} those of anchor \numCfgTuningBudgetToAnchorOneHundred{}, and the budgets
\numCfgBudgetsHostsItemFive{} and \numCfgBudgetsHostsItemSix{} those of anchor
\numCfgTuningBudgetToAnchorFourHundred{}.

\begin{table}[htbp]
\caption{Search space of the tuning, identical for S, both injection routes of D and P1
(\texttt{configs/config.yaml}, \texttt{tuning.search\_space}).}
\label{tab:space}
\centering\small
\begin{tabular}{@{}lll@{}}
\toprule
Setting & Distribution & Range or choices \\
\midrule
Learning rate & log-uniform & \numCfgTuningSearchSpaceLearningRateLow{} to \numCfgTuningSearchSpaceLearningRateHigh{} \\
Weight decay & log-uniform & \numCfgTuningSearchSpaceWeightDecayLow{} to \numCfgTuningSearchSpaceWeightDecayHigh{} \\
Hidden width & categorical & \numCfgTuningSearchSpaceHiddenWidthChoicesItemOne{}, \numCfgTuningSearchSpaceHiddenWidthChoicesItemTwo{} \\
MEGNet blocks & categorical & \numCfgTuningSearchSpaceMegnetBlocksChoicesItemOne{}, \numCfgTuningSearchSpaceMegnetBlocksChoicesItemTwo{}, \numCfgTuningSearchSpaceMegnetBlocksChoicesItemThree{} \\
Dropout & uniform & \numCfgTuningSearchSpaceDropoutLow{} to \numCfgTuningSearchSpaceDropoutHigh{} \\
Batch size (sites) & categorical & \numCfgTuningSearchSpaceBatchSizeChoicesItemOne{}, \numCfgTuningSearchSpaceBatchSizeChoicesItemTwo{}, \numCfgTuningSearchSpaceBatchSizeChoicesItemThree{} \\
Readout width & categorical & \numCfgTuningSearchSpaceReadoutMlpWidthChoicesItemOne{}, \numCfgTuningSearchSpaceReadoutMlpWidthChoicesItemTwo{} \\
Pooling of the atoms & categorical & Set2Set, mean \\
\bottomrule
\end{tabular}
\end{table}

The originally registered network (v1) was tuned under the frozen plan, and its values are those of the tag
\tagtuned{} (Table~\ref{tab:tuned_v1}). The two injection routes of D were tuned separately, and the rule logged
before tuning kept the route with the lower mean, over the anchors, of the best validation error: D-state, with
\numTuningSummaryDecisionMeansEVDState{}~eV against \numTuningSummaryDecisionMeansEVDLate{}~eV for D-late, a
difference of \numTuningSummaryDecisionDifferenceEV{}~eV. The two routes are indistinguishable on this objective,
and the selection is not evidence that one route encodes the electronic structure better. The primary network
(v2) differs from v1 only in its initial weights (He-uniform initialisation and zero biases); it was adopted after
test-set diagnostics of v1 and is post hoc. S and D of v2 were retuned with the same protocol, and test hosts were
never loaded in its design screen or its tuning (Table~\ref{tab:tuned_v2}). The descriptor predictor of v2, P1-v2,
was not retuned: it uses the v1 values of P1 with the v2 initialisation. Across both tables, the settings chosen at
the smallest anchor rest on very few validation hosts and should be read as one reasonable setting rather than
an optimum.

\begin{table}[htbp]
\caption{Tuned values of the originally registered network (v1; \texttt{configs/tuned\_v1.yaml}, tag \tagtuned{}),
with the best validation MAE (eV for S and D; the mean absolute error over the standardised descriptors,
dimensionless, for P1) of each study (\texttt{results/tuning\_summary.json}). D is the selected D-state route;
D-late is reported for completeness.}
\label{tab:tuned_v1}
\centering\footnotesize\tighttabsb
\begin{adjustbox}{max width=\linewidth}
% Generated by paper/si/tools/make_si_tables.py from configs/tuned_v1.yaml, results/tuning_summary.json. Do not edit.
\begin{tabular}{@{}llllllllllr@{}}
\toprule
Model & Anchor & Learning rate & Weight decay & Width & Blocks & Dropout & Batch & Readout & Pooling & Val.\ MAE \\
\midrule
S & \numCfgTuningAnchorsItemOne{} & \numTunedVOneModelsSFiftyLearningRateSig{} & \numTunedVOneModelsSFiftyWeightDecaySig{} & \numTunedVOneModelsSFiftyHiddenWidth{} & \numTunedVOneModelsSFiftyMegnetBlocks{} & \numTunedVOneModelsSFiftyDropoutSig{} & \numTunedVOneModelsSFiftyBatchSize{} & \numTunedVOneModelsSFiftyReadoutMlpWidth{} & set2set & \numTuningSummaryStudiesSAFiftyBestValue{} \\
 & \numCfgTuningAnchorsItemTwo{} & \numTunedVOneModelsSTwoHundredLearningRateSig{} & \numTunedVOneModelsSTwoHundredWeightDecaySig{} & \numTunedVOneModelsSTwoHundredHiddenWidth{} & \numTunedVOneModelsSTwoHundredMegnetBlocks{} & \numTunedVOneModelsSTwoHundredDropoutSig{} & \numTunedVOneModelsSTwoHundredBatchSize{} & \numTunedVOneModelsSTwoHundredReadoutMlpWidth{} & set2set & \numTuningSummaryStudiesSATwoHundredBestValue{} \\
 & \numCfgTuningAnchorsItemThree{} & \numTunedVOneModelsSSixHundredFiftyFourLearningRateSig{} & \numTunedVOneModelsSSixHundredFiftyFourWeightDecaySig{} & \numTunedVOneModelsSSixHundredFiftyFourHiddenWidth{} & \numTunedVOneModelsSSixHundredFiftyFourMegnetBlocks{} & \numTunedVOneModelsSSixHundredFiftyFourDropoutSig{} & \numTunedVOneModelsSSixHundredFiftyFourBatchSize{} & \numTunedVOneModelsSSixHundredFiftyFourReadoutMlpWidth{} & mean & \numTuningSummaryStudiesSASixHundredFiftyFourBestValue{} \\
\addlinespace
D (D-state) & \numCfgTuningAnchorsItemOne{} & \numTunedVOneModelsDStateFiftyLearningRateSig{} & \numTunedVOneModelsDStateFiftyWeightDecaySig{} & \numTunedVOneModelsDStateFiftyHiddenWidth{} & \numTunedVOneModelsDStateFiftyMegnetBlocks{} & \numTunedVOneModelsDStateFiftyDropoutSig{} & \numTunedVOneModelsDStateFiftyBatchSize{} & \numTunedVOneModelsDStateFiftyReadoutMlpWidth{} & set2set & \numTuningSummaryStudiesDStateAFiftyBestValue{} \\
 & \numCfgTuningAnchorsItemTwo{} & \numTunedVOneModelsDStateTwoHundredLearningRateSig{} & \numTunedVOneModelsDStateTwoHundredWeightDecaySig{} & \numTunedVOneModelsDStateTwoHundredHiddenWidth{} & \numTunedVOneModelsDStateTwoHundredMegnetBlocks{} & \numTunedVOneModelsDStateTwoHundredDropoutSig{} & \numTunedVOneModelsDStateTwoHundredBatchSize{} & \numTunedVOneModelsDStateTwoHundredReadoutMlpWidth{} & set2set & \numTuningSummaryStudiesDStateATwoHundredBestValue{} \\
 & \numCfgTuningAnchorsItemThree{} & \numTunedVOneModelsDStateSixHundredFiftyFourLearningRateSig{} & \numTunedVOneModelsDStateSixHundredFiftyFourWeightDecaySig{} & \numTunedVOneModelsDStateSixHundredFiftyFourHiddenWidth{} & \numTunedVOneModelsDStateSixHundredFiftyFourMegnetBlocks{} & \numTunedVOneModelsDStateSixHundredFiftyFourDropoutSig{} & \numTunedVOneModelsDStateSixHundredFiftyFourBatchSize{} & \numTunedVOneModelsDStateSixHundredFiftyFourReadoutMlpWidth{} & mean & \numTuningSummaryStudiesDStateASixHundredFiftyFourBestValue{} \\
\addlinespace
D-late & \numCfgTuningAnchorsItemOne{} & \numTunedVOneModelsDLateFiftyLearningRateSig{} & \numTunedVOneModelsDLateFiftyWeightDecaySig{} & \numTunedVOneModelsDLateFiftyHiddenWidth{} & \numTunedVOneModelsDLateFiftyMegnetBlocks{} & \numTunedVOneModelsDLateFiftyDropoutSig{} & \numTunedVOneModelsDLateFiftyBatchSize{} & \numTunedVOneModelsDLateFiftyReadoutMlpWidth{} & mean & \numTuningSummaryStudiesDLateAFiftyBestValue{} \\
 & \numCfgTuningAnchorsItemTwo{} & \numTunedVOneModelsDLateTwoHundredLearningRateSig{} & \numTunedVOneModelsDLateTwoHundredWeightDecaySig{} & \numTunedVOneModelsDLateTwoHundredHiddenWidth{} & \numTunedVOneModelsDLateTwoHundredMegnetBlocks{} & \numTunedVOneModelsDLateTwoHundredDropoutSig{} & \numTunedVOneModelsDLateTwoHundredBatchSize{} & \numTunedVOneModelsDLateTwoHundredReadoutMlpWidth{} & set2set & \numTuningSummaryStudiesDLateATwoHundredBestValue{} \\
 & \numCfgTuningAnchorsItemThree{} & \numTunedVOneModelsDLateSixHundredFiftyFourLearningRateSig{} & \numTunedVOneModelsDLateSixHundredFiftyFourWeightDecaySig{} & \numTunedVOneModelsDLateSixHundredFiftyFourHiddenWidth{} & \numTunedVOneModelsDLateSixHundredFiftyFourMegnetBlocks{} & \numTunedVOneModelsDLateSixHundredFiftyFourDropoutSig{} & \numTunedVOneModelsDLateSixHundredFiftyFourBatchSize{} & \numTunedVOneModelsDLateSixHundredFiftyFourReadoutMlpWidth{} & mean & \numTuningSummaryStudiesDLateASixHundredFiftyFourBestValue{} \\
\addlinespace
P1 & \numCfgTuningAnchorsItemOne{} & \numTunedVOneModelsPOneFiftyLearningRateSig{} & \numTunedVOneModelsPOneFiftyWeightDecaySig{} & \numTunedVOneModelsPOneFiftyHiddenWidth{} & \numTunedVOneModelsPOneFiftyMegnetBlocks{} & \numTunedVOneModelsPOneFiftyDropoutSig{} & \numTunedVOneModelsPOneFiftyBatchSize{} & \numTunedVOneModelsPOneFiftyReadoutMlpWidth{} & set2set & \numTuningSummaryStudiesPOneAFiftyBestValue{} \\
 & \numCfgTuningAnchorsItemTwo{} & \numTunedVOneModelsPOneTwoHundredLearningRateSig{} & \numTunedVOneModelsPOneTwoHundredWeightDecaySig{} & \numTunedVOneModelsPOneTwoHundredHiddenWidth{} & \numTunedVOneModelsPOneTwoHundredMegnetBlocks{} & \numTunedVOneModelsPOneTwoHundredDropoutSig{} & \numTunedVOneModelsPOneTwoHundredBatchSize{} & \numTunedVOneModelsPOneTwoHundredReadoutMlpWidth{} & mean & \numTuningSummaryStudiesPOneATwoHundredBestValue{} \\
 & \numCfgTuningAnchorsItemThree{} & \numTunedVOneModelsPOneSixHundredFiftyFourLearningRateSig{} & \numTunedVOneModelsPOneSixHundredFiftyFourWeightDecaySig{} & \numTunedVOneModelsPOneSixHundredFiftyFourHiddenWidth{} & \numTunedVOneModelsPOneSixHundredFiftyFourMegnetBlocks{} & \numTunedVOneModelsPOneSixHundredFiftyFourDropoutSig{} & \numTunedVOneModelsPOneSixHundredFiftyFourBatchSize{} & \numTunedVOneModelsPOneSixHundredFiftyFourReadoutMlpWidth{} & mean & \numTuningSummaryStudiesPOneASixHundredFiftyFourBestValue{} \\
\bottomrule
\end{tabular}

\end{adjustbox}
\end{table}

\begin{table}[htbp]
\caption{Tuned values of the primary network (v2, post hoc; \texttt{configs/tuned\_v2.yaml}). Same protocol and
search space as Table~\ref{tab:tuned_v1}; P1-v2 reuses the P1 values of Table~\ref{tab:tuned_v1}.}
\label{tab:tuned_v2}
\centering\footnotesize\tighttabsb
\begin{adjustbox}{max width=\linewidth}
% Generated by paper/si/tools/make_si_tables.py from configs/tuned_v2.yaml. Do not edit.
\begin{tabular}{@{}llllllllll@{}}
\toprule
Model & Anchor & Learning rate & Weight decay & Width & Blocks & Dropout & Batch & Readout & Pooling \\
\midrule
S & \numCfgTuningAnchorsItemOne{} & \numTunedVTwoModelsSFiftyLearningRateSig{} & \numTunedVTwoModelsSFiftyWeightDecaySig{} & \numTunedVTwoModelsSFiftyHiddenWidth{} & \numTunedVTwoModelsSFiftyMegnetBlocks{} & \numTunedVTwoModelsSFiftyDropoutSig{} & \numTunedVTwoModelsSFiftyBatchSize{} & \numTunedVTwoModelsSFiftyReadoutMlpWidth{} & mean \\
 & \numCfgTuningAnchorsItemTwo{} & \numTunedVTwoModelsSTwoHundredLearningRateSig{} & \numTunedVTwoModelsSTwoHundredWeightDecaySig{} & \numTunedVTwoModelsSTwoHundredHiddenWidth{} & \numTunedVTwoModelsSTwoHundredMegnetBlocks{} & \numTunedVTwoModelsSTwoHundredDropoutSig{} & \numTunedVTwoModelsSTwoHundredBatchSize{} & \numTunedVTwoModelsSTwoHundredReadoutMlpWidth{} & mean \\
 & \numCfgTuningAnchorsItemThree{} & \numTunedVTwoModelsSSixHundredFiftyFourLearningRateSig{} & \numTunedVTwoModelsSSixHundredFiftyFourWeightDecaySig{} & \numTunedVTwoModelsSSixHundredFiftyFourHiddenWidth{} & \numTunedVTwoModelsSSixHundredFiftyFourMegnetBlocks{} & \numTunedVTwoModelsSSixHundredFiftyFourDropoutSig{} & \numTunedVTwoModelsSSixHundredFiftyFourBatchSize{} & \numTunedVTwoModelsSSixHundredFiftyFourReadoutMlpWidth{} & set2set \\
\addlinespace
D (D-state) & \numCfgTuningAnchorsItemOne{} & \numTunedVTwoModelsDStateFiftyLearningRateSig{} & \numTunedVTwoModelsDStateFiftyWeightDecaySig{} & \numTunedVTwoModelsDStateFiftyHiddenWidth{} & \numTunedVTwoModelsDStateFiftyMegnetBlocks{} & \numTunedVTwoModelsDStateFiftyDropoutSig{} & \numTunedVTwoModelsDStateFiftyBatchSize{} & \numTunedVTwoModelsDStateFiftyReadoutMlpWidth{} & mean \\
 & \numCfgTuningAnchorsItemTwo{} & \numTunedVTwoModelsDStateTwoHundredLearningRateSig{} & \numTunedVTwoModelsDStateTwoHundredWeightDecaySig{} & \numTunedVTwoModelsDStateTwoHundredHiddenWidth{} & \numTunedVTwoModelsDStateTwoHundredMegnetBlocks{} & \numTunedVTwoModelsDStateTwoHundredDropoutSig{} & \numTunedVTwoModelsDStateTwoHundredBatchSize{} & \numTunedVTwoModelsDStateTwoHundredReadoutMlpWidth{} & mean \\
 & \numCfgTuningAnchorsItemThree{} & \numTunedVTwoModelsDStateSixHundredFiftyFourLearningRateSig{} & \numTunedVTwoModelsDStateSixHundredFiftyFourWeightDecaySig{} & \numTunedVTwoModelsDStateSixHundredFiftyFourHiddenWidth{} & \numTunedVTwoModelsDStateSixHundredFiftyFourMegnetBlocks{} & \numTunedVTwoModelsDStateSixHundredFiftyFourDropoutSig{} & \numTunedVTwoModelsDStateSixHundredFiftyFourBatchSize{} & \numTunedVTwoModelsDStateSixHundredFiftyFourReadoutMlpWidth{} & mean \\
\bottomrule
\end{tabular}

\end{adjustbox}
\end{table}

\subsection{The cap on training length and its measured effects}

The cap of \numCfgTrainingMaxEpochs{} passes over the training hosts and the patience of
\numCfgTrainingEarlyStoppingPatience{} passes were set by a rule logged before the calibration pilot ran: the cap
is a round number of passes above the slowest convergence seen in the pilot, with a safety factor, and the
patience follows from the cap (\texttt{docs/deviations.md}, the \sessioncap{} rows; \texttt{docs/methods\_notes.md}). The
registered check of the cap retrained S and D with a cap of \numGOneCapSensitivityCapsItemTwo{} passes at
\numGOneCapSensitivityBudget{} hosts on \numGOneCapSensitivityBootstrapNResamples{} resamples and all seeds. All
\numGOneDiagnosticsEpochCapPairsNPairs{} pairs stopped at the same pass and gave identical predictions
(\numGOneDiagnosticsEpochCapPairsNIdentical{} of \numGOneDiagnosticsEpochCapPairsNPairs{}), so the advantage of the
descriptors was \numGOneCapSensitivityATwoHundredMean{}~eV under both caps
\ci{\numGOneCapSensitivityATwoHundredCiLo{}}{\numGOneCapSensitivityATwoHundredCiHi{}} and the primary result was
cap-insensitive by the registered criterion (registered v1; \texttt{results/g1\_cap\_sensitivity.json}).

The extensions show that the cap binds more often than those runs suggested, and mostly for the structure-only
network (Table~\ref{tab:capped}). Among the precision runs of v1 at \numCfgBudgetsHostsItemSix{} hosts,
\numPrecisionABFourHundredBSixHundredFiftyFourDescriptiveCappedPrecisionSixHundredFiftyFourVOneSNCapped{} of
\numPrecisionABFourHundredBSixHundredFiftyFourDescriptiveCappedPrecisionSixHundredFiftyFourVOneSNRuns{} S runs
reached the cap, against
\numPrecisionABFourHundredBSixHundredFiftyFourDescriptiveCappedPrecisionSixHundredFiftyFourVOneDNCapped{} D runs;
for v2 the counts were
\numPrecisionABFourHundredBSixHundredFiftyFourDescriptiveCappedPrecisionSixHundredFiftyFourVTwoSNCapped{} and
\numPrecisionABFourHundredBSixHundredFiftyFourDescriptiveCappedPrecisionSixHundredFiftyFourVTwoDNCapped{}. In every
group of S runs with capped members, the capped runs had the lower mean error, so the structure-only network may
still have been improving when it was stopped at the largest training sets. Because capped and uncapped runs differ
in resample and fold as well, Table~\ref{tab:capped} is descriptive and carries no test. The post hoc rerun of every
capped primary run with a longer cap is reported in Section~\ref{si:sensitivities}: it changes the advantage at
\numCfgBudgetsHostsItemSix{} hosts by
\numSensCappedPostHocPerBudgetSixHundredFiftyFourChangeTwinMinusPrimaryAMean{}~eV on v1, and no such reruns exist
for v2.

\begin{table}[htbp]
\caption{Runs that stopped at the cap of \numCfgTrainingMaxEpochs{} passes (extension, descriptive, no test), and
the mean single-run site MAE (eV) of capped and uncapped runs. Precision: the pooled runs at the two largest
budgets; grouped CV: all folds at each $K$. Records: \recprecision{} and \reccv{}, key
\texttt{payload.descriptive.capped}.}
\label{tab:capped}
\centering\scriptsize
% Generated by paper/si/tools/make_si_tables.py from results/precision_a_b400_b654.json, results/cv_extended_budget.json (payload.descriptive.capped). Do not edit.
\begin{tabular}{@{}llllrrrr@{}}
\toprule
Design & Budget or folds & Network & Model & Runs & Stopped at cap & MAE, capped & MAE, not capped \\
\midrule
precision & \numCfgBudgetsHostsItemFive{} & v2 & S & \numPrecisionABFourHundredBSixHundredFiftyFourDescriptiveCappedPrecisionFourHundredVTwoSNRuns{} & \numPrecisionABFourHundredBSixHundredFiftyFourDescriptiveCappedPrecisionFourHundredVTwoSNCapped{} & \numPrecisionABFourHundredBSixHundredFiftyFourDescriptiveCappedPrecisionFourHundredVTwoSMeanSingleRunMaeCapped{} & \numPrecisionABFourHundredBSixHundredFiftyFourDescriptiveCappedPrecisionFourHundredVTwoSMeanSingleRunMaeUncapped{} \\
precision &  & v2 & D & \numPrecisionABFourHundredBSixHundredFiftyFourDescriptiveCappedPrecisionFourHundredVTwoDNRuns{} & \numPrecisionABFourHundredBSixHundredFiftyFourDescriptiveCappedPrecisionFourHundredVTwoDNCapped{} & -- & \numPrecisionABFourHundredBSixHundredFiftyFourDescriptiveCappedPrecisionFourHundredVTwoDMeanSingleRunMaeUncapped{} \\
precision &  & v1 & S & \numPrecisionABFourHundredBSixHundredFiftyFourDescriptiveCappedPrecisionFourHundredVOneSNRuns{} & \numPrecisionABFourHundredBSixHundredFiftyFourDescriptiveCappedPrecisionFourHundredVOneSNCapped{} & \numPrecisionABFourHundredBSixHundredFiftyFourDescriptiveCappedPrecisionFourHundredVOneSMeanSingleRunMaeCapped{} & \numPrecisionABFourHundredBSixHundredFiftyFourDescriptiveCappedPrecisionFourHundredVOneSMeanSingleRunMaeUncapped{} \\
precision &  & v1 & D & \numPrecisionABFourHundredBSixHundredFiftyFourDescriptiveCappedPrecisionFourHundredVOneDNRuns{} & \numPrecisionABFourHundredBSixHundredFiftyFourDescriptiveCappedPrecisionFourHundredVOneDNCapped{} & \numPrecisionABFourHundredBSixHundredFiftyFourDescriptiveCappedPrecisionFourHundredVOneDMeanSingleRunMaeCapped{} & \numPrecisionABFourHundredBSixHundredFiftyFourDescriptiveCappedPrecisionFourHundredVOneDMeanSingleRunMaeUncapped{} \\
\addlinespace
precision & \numCfgBudgetsHostsItemSix{} & v2 & S & \numPrecisionABFourHundredBSixHundredFiftyFourDescriptiveCappedPrecisionSixHundredFiftyFourVTwoSNRuns{} & \numPrecisionABFourHundredBSixHundredFiftyFourDescriptiveCappedPrecisionSixHundredFiftyFourVTwoSNCapped{} & \numPrecisionABFourHundredBSixHundredFiftyFourDescriptiveCappedPrecisionSixHundredFiftyFourVTwoSMeanSingleRunMaeCapped{} & \numPrecisionABFourHundredBSixHundredFiftyFourDescriptiveCappedPrecisionSixHundredFiftyFourVTwoSMeanSingleRunMaeUncapped{} \\
precision &  & v2 & D & \numPrecisionABFourHundredBSixHundredFiftyFourDescriptiveCappedPrecisionSixHundredFiftyFourVTwoDNRuns{} & \numPrecisionABFourHundredBSixHundredFiftyFourDescriptiveCappedPrecisionSixHundredFiftyFourVTwoDNCapped{} & \numPrecisionABFourHundredBSixHundredFiftyFourDescriptiveCappedPrecisionSixHundredFiftyFourVTwoDMeanSingleRunMaeCapped{} & \numPrecisionABFourHundredBSixHundredFiftyFourDescriptiveCappedPrecisionSixHundredFiftyFourVTwoDMeanSingleRunMaeUncapped{} \\
precision &  & v1 & S & \numPrecisionABFourHundredBSixHundredFiftyFourDescriptiveCappedPrecisionSixHundredFiftyFourVOneSNRuns{} & \numPrecisionABFourHundredBSixHundredFiftyFourDescriptiveCappedPrecisionSixHundredFiftyFourVOneSNCapped{} & \numPrecisionABFourHundredBSixHundredFiftyFourDescriptiveCappedPrecisionSixHundredFiftyFourVOneSMeanSingleRunMaeCapped{} & \numPrecisionABFourHundredBSixHundredFiftyFourDescriptiveCappedPrecisionSixHundredFiftyFourVOneSMeanSingleRunMaeUncapped{} \\
precision &  & v1 & D & \numPrecisionABFourHundredBSixHundredFiftyFourDescriptiveCappedPrecisionSixHundredFiftyFourVOneDNRuns{} & \numPrecisionABFourHundredBSixHundredFiftyFourDescriptiveCappedPrecisionSixHundredFiftyFourVOneDNCapped{} & \numPrecisionABFourHundredBSixHundredFiftyFourDescriptiveCappedPrecisionSixHundredFiftyFourVOneDMeanSingleRunMaeCapped{} & \numPrecisionABFourHundredBSixHundredFiftyFourDescriptiveCappedPrecisionSixHundredFiftyFourVOneDMeanSingleRunMaeUncapped{} \\
\addlinespace
grouped CV & $K = \numCvExtendedBudgetDescriptiveCappedCvTenVTwoSPoint{}$ & v2 & S & \numCvExtendedBudgetDescriptiveCappedCvTenVTwoSNRuns{} & \numCvExtendedBudgetDescriptiveCappedCvTenVTwoSNCapped{} & \numCvExtendedBudgetDescriptiveCappedCvTenVTwoSMeanSingleRunMaeCapped{} & \numCvExtendedBudgetDescriptiveCappedCvTenVTwoSMeanSingleRunMaeUncapped{} \\
grouped CV &  & v2 & D & \numCvExtendedBudgetDescriptiveCappedCvTenVTwoDNRuns{} & \numCvExtendedBudgetDescriptiveCappedCvTenVTwoDNCapped{} & \numCvExtendedBudgetDescriptiveCappedCvTenVTwoDMeanSingleRunMaeCapped{} & \numCvExtendedBudgetDescriptiveCappedCvTenVTwoDMeanSingleRunMaeUncapped{} \\
grouped CV &  & v1 & S & \numCvExtendedBudgetDescriptiveCappedCvTenVOneSNRuns{} & \numCvExtendedBudgetDescriptiveCappedCvTenVOneSNCapped{} & \numCvExtendedBudgetDescriptiveCappedCvTenVOneSMeanSingleRunMaeCapped{} & \numCvExtendedBudgetDescriptiveCappedCvTenVOneSMeanSingleRunMaeUncapped{} \\
grouped CV &  & v1 & D & \numCvExtendedBudgetDescriptiveCappedCvTenVOneDNRuns{} & \numCvExtendedBudgetDescriptiveCappedCvTenVOneDNCapped{} & \numCvExtendedBudgetDescriptiveCappedCvTenVOneDMeanSingleRunMaeCapped{} & \numCvExtendedBudgetDescriptiveCappedCvTenVOneDMeanSingleRunMaeUncapped{} \\
\addlinespace
grouped CV & $K = \numCvExtendedBudgetDescriptiveCappedCvTwentyVTwoSPoint{}$ & v2 & S & \numCvExtendedBudgetDescriptiveCappedCvTwentyVTwoSNRuns{} & \numCvExtendedBudgetDescriptiveCappedCvTwentyVTwoSNCapped{} & \numCvExtendedBudgetDescriptiveCappedCvTwentyVTwoSMeanSingleRunMaeCapped{} & \numCvExtendedBudgetDescriptiveCappedCvTwentyVTwoSMeanSingleRunMaeUncapped{} \\
grouped CV &  & v2 & D & \numCvExtendedBudgetDescriptiveCappedCvTwentyVTwoDNRuns{} & \numCvExtendedBudgetDescriptiveCappedCvTwentyVTwoDNCapped{} & \numCvExtendedBudgetDescriptiveCappedCvTwentyVTwoDMeanSingleRunMaeCapped{} & \numCvExtendedBudgetDescriptiveCappedCvTwentyVTwoDMeanSingleRunMaeUncapped{} \\
grouped CV &  & v1 & S & \numCvExtendedBudgetDescriptiveCappedCvTwentyVOneSNRuns{} & \numCvExtendedBudgetDescriptiveCappedCvTwentyVOneSNCapped{} & \numCvExtendedBudgetDescriptiveCappedCvTwentyVOneSMeanSingleRunMaeCapped{} & \numCvExtendedBudgetDescriptiveCappedCvTwentyVOneSMeanSingleRunMaeUncapped{} \\
grouped CV &  & v1 & D & \numCvExtendedBudgetDescriptiveCappedCvTwentyVOneDNRuns{} & \numCvExtendedBudgetDescriptiveCappedCvTwentyVOneDNCapped{} & \numCvExtendedBudgetDescriptiveCappedCvTwentyVOneDMeanSingleRunMaeCapped{} & \numCvExtendedBudgetDescriptiveCappedCvTwentyVOneDMeanSingleRunMaeUncapped{} \\
\bottomrule
\end{tabular}

\end{table}

\subsection{Hardware}

The registered tuning and the registered sweep of the originally registered network (v1) ran on one
\gpulfortys{}, with several training processes sharing the card; a process started only when the memory of the
running jobs left a fixed margin (\texttt{docs/methods\_notes.md}). These settings change wall time only, not any
setting, split, seed or stopping rule. The run records of the original sweep, and those of the collaborator's v2
tuning and production, name a CUDA device but not the GPU model. The extension training runs (precision, grouped
cross-validation, high-moment, CGCNN and capped twins) ran on the GPUs of the National Research Platform, one
training process per GPU, with PyTorch \vertorch{} and matgl \vermatgl{}; their records name the GPU model:
\gpuaten{}, \gputhreeoh{}, \gpulfour{} and \gpufourty{}. The inference stages of the latent probe and the geometry
evaluation also used an \gputitan{} (\texttt{docs/nrp/n3\_ingest.md}). Training on a GPU is not bitwise
reproducible across GPU models, because the summation order of the aggregation over neighbours differs;
Section~\ref{si:hardware} gives the repeatability study and the size of the effect on the predicted formation
energies.
